% ============================================================
% 3. System Implementation \& Tooling
% ============================================================
\section{System Implementation \& Tooling}
\label{sec:implementation}

\subsection{Data Generation: Synthetic Corruption Pipeline}
\label{sec:data_generation}

The synthetic data generator (\texttt{generate\_data.py}, \texttt{data\_generator/}) creates two datasets A and B with known \globalentity{} overlap.
The pipeline (\Cref{fig:datagen_pipeline}):

\begin{figure}[htbp]
	\centering
	\makebox[\textwidth][c]{
		\begin{tikzpicture}[
				node distance=0.8cm,
				every node/.style={font=\small},
				stage/.style={
						draw,
						rectangle,
						rounded corners=4pt,
						minimum width=11cm,
						minimum height=1.1cm,
						align=center,
						thick,
						inner xsep=10pt,
						inner ysep=6pt
					},
				arrow/.style={->, thick, shorten >=2pt, shorten <=2pt}
			]
		% Pipeline stages - vertical stack
		\node[stage, fill=blue!10] (step1) {
			\textbf{1. EntitySpaceManager} \quad Creates $N$ global entities with overlap rate $\rho$
		};
		\node[stage, fill=green!10, below=of step1] (step2) {
			\textbf{2. Source Mappings $M_A, M_B$} \quad Assign to datasets (shared/unique split)
		};
		\node[stage, fill=yellow!10, below=of step2] (step3) {
			\textbf{3. DataPipelineOrchestrator} \quad Expands each entity $\to$ multiple records (cardinality per entity)
		};
		\node[stage, fill=orange!10, below=of step3] (step4) {
			\textbf{4. Asymmetric Corruption} \quad Independent noise per source (typos, missing fields, date drift)
		};
		\node[stage, fill=red!10, below=of step4] (step5) {
			\textbf{5. Latent IDs Preserved} \quad \texttt{OVER\_*}, \texttt{UNI\_A\_*}, \texttt{UNIQ\_B\_*} for ground-truth evaluation
		};

		% Stacked arrows
		\foreach \src/\dst in {step1/step2, step2/step3, step3/step4, step4/step5} {
				\draw[arrow] (\src.south) -- (\dst.north);
			}

		% Side labels showing data flow transformation (Fixed coordinate syntax)
		\node[anchor=east, font=\tiny\ttfamily, text=gray] at ([xshift=-0.2cm]step1.west) {Pool of $N$ entities};
		\node[anchor=east, font=\tiny\ttfamily, text=gray] at ([xshift=-0.2cm]step2.west) {Split: $R_A$, $R_B$};
		\node[anchor=east, font=\tiny\ttfamily, text=gray] at ([xshift=-0.2cm]step3.west) {Multi-record per entity};
		\node[anchor=east, font=\tiny\ttfamily, text=gray] at ([xshift=-0.2cm]step4.west) {Source-specific noise};
		\node[anchor=east, font=\tiny\ttfamily, text=gray] at ([xshift=-0.2cm]step5.west) {Evaluation keys};
	\end{tikzpicture}
	}
	\caption{Synthetic Data Generation Pipeline.
		Vertical stages show the transformation from entity pool $\to$ split datasets $\to$ multi-record expansion $\to$ independent corruption $\to$ ground-truth keys.
		Each stage is a separable, testable component in the codebase.}
	\label{fig:datagen_pipeline}
\end{figure}

The pipeline was designed to easily support customisations at various levels (config, data type, data distributions).

\subsubsection{Configuration (YAML)}
\label{sec:datagen_config}
Key parameters (\texttt{data\_generator/config.yaml}):
\begin{itemize}
	\item \texttt{dataset\_size}: Total global entities $N$.
	\item \texttt{overlap\_percentage}: Fraction $\rho$ of entities present in both datasets.
	\item \texttt{source\_a\_rows\_per\_entity}, \texttt{source\_b\_rows\_per\_entity}: Multiplicity (cardinality) per entity per source.
	\item \texttt{features}: List of $\{ \text{name, type, noise\_params, missing\_rate} \}$.
	      Types: \texttt{categorical} (classes, missing), \texttt{numerical} (Gaussian noise $\sigma$), \texttt{datetime} (Gaussian hours, missing).
	\item \texttt{overlap\_dropout\_frac}: Asymmetric dropout on overlapping portion of B.
\end{itemize}

\subsubsection{Noise Models}
\label{sec:noise_models}
\begin{itemize}
	\item \textbf{Categorical}: Uniform random generation out of \texttt{num\_classes}, missing with \texttt{missing\_rate}.
	\item \textbf{Numerical}: True value $+\ \mathcal{N}(0, \sigma^2)$; missing with \texttt{missing\_rate}.
	\item \textbf{Datetime}: True timestamp $+\ \mathcal{N}(0, \sigma_{\text{hours}}^2)$ hours; missing with \texttt{missing\_rate}.
\end{itemize}

\subsubsection{Latent Event IDs}
Every generated row receives a \texttt{latent\_event\_id}:
\begin{itemize}
	\item \texttt{OVER\_k}: Overlapping entities (same $k$ in A and B $\to$ true match).
	\item \texttt{UNI\_A\_k}: Unique to A.
	\item \texttt{UNIQ\_B\_k}: Unique to B.
\end{itemize}
These enable ground-truth evaluation without manual labelling.


\subsection{Profile Aggregation Engine}
\label{sec:agg_profile_engine}
The profile aggregation engine transforms raw record tables into entity-level profile vectors.
A simplified horizontal flow diagram illustrates the transformation:

\begin{figure}[htbp]
	\centering
	\begin{tikzpicture}[
			node distance=1.0cm,
			every node/.style={font=\small},
			tableTitle/.style={font=\small\bfseries},
			% Compact stealth arrows
			arrow/.style={->, >=stealth, thick, line width=1.2pt, color=gray!80, shorten >=6pt, shorten <=6pt}
		]

		% --- INPUT TABLE ---
		\matrix [matrix of nodes, nodes={align=left, inner sep=3pt}, row sep=1pt, column sep=4pt] (inputTable) {
			\rowcolor{tableHeader} \textbf{\texttt{unique\_id}} & \textbf{\texttt{feature\_1}} & \textbf{\texttt{feature\_2}} & \textbf{\texttt{feature\_3}} \\
			\texttt{e1} & \texttt{val\_a1} & \texttt{val\_b1} & \texttt{val\_c1} \\
			\rowcolor{tableRowEven} \texttt{e1} & \texttt{val\_a2} & \texttt{val\_b2} & \texttt{val\_c2} \\
			\texttt{e2} & \texttt{val\_a3} & \texttt{val\_b3} & \texttt{val\_c3} \\
		};
		% Crisp top/bottom borders
		\draw[line width=1.1pt] (inputTable.north west) -- (inputTable.north east);
		\draw[line width=0.6pt] (inputTable-1-1.south west) -- (inputTable-1-4.south east);
		\draw[line width=1.1pt] (inputTable.south west) -- (inputTable.south east);
		\node[tableTitle, above=0.05cm of inputTable] (inputTitle) {Input Table};

		% --- OUTPUT PROFILE TABLE ---
		\matrix [matrix of nodes, nodes={align=left, inner sep=3pt}, row sep=1pt, column sep=4pt, right=1.5cm of inputTable] (outputTable) {
			\rowcolor{tableHeader} \textbf{\texttt{unique\_id}} & \textbf{\texttt{prof\_1}} & \textbf{\texttt{prof\_2}} & \textbf{\texttt{prof\_3}} \\
			\texttt{e1} & \texttt{agg\_a} & \texttt{agg\_b} & \texttt{agg\_c} \\
			\rowcolor{tableRowEven} \texttt{e2} & \texttt{agg\_a'} & \texttt{agg\_b'} & \texttt{agg\_c'} \\
		};
		% Crisp top/bottom borders matching the input table exactly
		\draw[line width=1.1pt] (outputTable.north west) -- (outputTable.north east);
		\draw[line width=0.6pt] (outputTable-1-1.south west) -- (outputTable-1-4.south east);
		\draw[line width=1.1pt] (outputTable.south west) -- (outputTable.south east);
		\node[tableTitle, above=0.05cm of outputTable] (outputTitle) {Profile Table};

		% --- DIRECT FLOW ARROW WITH LABEL ---
		\draw[arrow] (inputTable.east) -- node[above, font=\scriptsize\bfseries, color=black!70] {Aggregate} (outputTable.west);

	\end{tikzpicture}
	\caption{Horizontal flow of the profile aggregation engine: raw records are grouped by entity identifier and transformed into fixed‑dimensional profile vectors.}
	\label{fig:profile_agg_flow}
\end{figure}

The full implementation resides in the \texttt{em/aggregation/} package. It uses a \texttt{ProcessorRegistry} to map feature types to \texttt{ColumnProcessor} subclasses, and an \texttt{Aggregator} that coordinates multiple processors, builds vocabularies, applies aggregation functions, and runs post‑processors.

\subsection{Splink Configuration Overview}
\label{sec:splink_config}
The Entity Matching pipeline uses the Splink library for probabilistic record linkage.
All synthetic datasets share identical comparison logic; only dataset paths and resource parameters vary.
Notebooks with real datasets reuse the same structure with domain‑specific field mappings.

Supported comparison functions include:
\begin{itemize}
	\item Basic Splink levels: \texttt{NullLevel}, \texttt{ExactMatchLevel}, \texttt{ElseLevel}.
	\item Binned levels: \texttt{PercentageDifferenceLevel}, \texttt{WithinNDaysLevel}, \texttt{JaroWinklerLevel}.
	\item Custom comparisons via SQL expressions (e.g., for set similarities, range overlap).
\end{itemize}

The core workflow follows the Splink transaction demo pattern:
\begin{enumerate}
	\item Define comparison levels for each field.
	\item Specify blocking rules as SQL predicates.
	\item Train the model via expectation‑maximisation (EM) using a fixed schedule of blocking rules.
	\item Generate predictions and evaluate against known latent IDs.
\end{enumerate}

Full code listings are in the respective notebooks mentioned by the report.

\subsection{Blocking Rules Design}
\label{sec:blocking_rules}
Blocking rules are supplied as an array of SQL predicates to the argument \texttt{blocking\_rules\_to\_generate\_predictions}.
Each rule targets a different matching signal (e.g., categorical equality, numerical proximity, temporal proximity) and their union bounds recall while reducing the comparison space.
Example predicates (for synthetic data) are of the form:
\begin{itemize}
	\item \texttt{l.categorical = r.categorical OR l.categorical IS NULL OR r.categorical IS NULL}
	\item \texttt{strftime(l.datetime\_noisy, '\%Y\%m') = strftime(r.datetime\_noisy, '\%Y\%m') OR l.datetime\_noisy IS NULL OR r.datetime\_noisy IS NULL}
	\item \texttt{(ABS(l.numerical\_noisy - r.numerical\_noisy) / NULLIF(CASE WHEN l.numerical\_noisy > r.numerical\_noisy THEN l.numerical\_noisy ELSE r.numerical\_noisy END, 0)) < 0.25 OR l.numerical\_noisy IS NULL OR r.numerical\_noisy IS NULL}
\end{itemize}

Full code listings are in the respective notebooks mentioned by the report.

\subsection{EM Training Schedule}
\label{sec:em_training}
Parameter estimation proceeds via a fixed sequence of expectation‑maximisation steps, each using a different blocking rule to estimate m‑probabilities for the remaining comparisons.
The schedule begins with estimating the global prior from strict rules, then estimates u‑probabilities via random sampling, followed by EM rounds that alternate blocking columns to identify all m‑probabilities.

Full code listings are in the respective notebooks mentioned by the report.
